% Generated draft: requires mathematical and source review.
\documentclass[11pt]{article}
\usepackage[margin=25mm]{geometry}
\usepackage{fontspec}
\setmainfont{Latin Modern Roman}
\usepackage{amsmath,amssymb,mathtools,unicode-math}
\setmathfont{Latin Modern Math}
\usepackage{enumitem,needspace,xurl,hyperref}
\setlength{\parindent}{0pt}
\setlength{\parskip}{5pt}
\setlength{\emergencystretch}{3em}
\newcommand{\N}{\mathbb{N}}
\newcommand{\Z}{\mathbb{Z}}
\newcommand{\Q}{\mathbb{Q}}
\newcommand{\R}{\mathbb{R}}
\newcommand{\C}{\mathbb{C}}
\newcommand{\F}{\mathbb{F}}
\newcommand{\A}{\mathbb{A}}
\newcommand{\PP}{\mathbb{P}}
\newcommand{\E}{\mathbb{E}}
\newcommand{\eps}{\varepsilon}
\newcommand{\dd}{\,\mathrm{d}}
\begin{document}
\section*{UnsolvedMath ID: 11000214}
\textbf{Source problem code:} catalogue entry
\subsection{Definitions and mathematical statement}
In the customary finite-area setting, a translation surface is a surface with flat coordinate changes $z\mapsto z+c$ away from finitely many singularities, equivalently a suitable pair $(X,\omega)$ with an abelian differential. An orientation-preserving affine self-map has a constant linear derivative, and its derivatives form the Veech group
\[
 \Gamma(X,\omega)=\{D a:a\in\operatorname{Aff}^+(X,\omega)\}\le\mathrm{SL}_2(\R).
\]
For a quadratic differential, the flat changes can be $z\mapsto\pm z+c$; the derivative is then naturally defined modulo sign, giving the projective Veech group in $\mathrm{PSL}_2(\R)$. The area, singularity/pole, and sign conventions need to follow the intended surface class.

\textbf{Problem 7 (algorithm for Veech groups).} Is there an algorithm that determines the full Veech group of a general translation surface or quadratic differential? A precise algorithmic formulation requires an effective encoding of the input, such as polygonal gluing data with specified effectively represented periods, and a stated kind of exact group description as output. These are not supplied in the excerpt. Arbitrary real coordinates cannot be treated as finite exact input without a computation convention. Likewise, determining a group could mean a membership decision procedure, a generator description, or a fundamental-domain construction; these are not identical tasks. No general finite-generation hypothesis or finite presentation is imposed by the source.

The source highlights square-tiled surfaces as an interesting class of Veech surfaces: they are obtained by gluing finitely many equal-sized squares. This is a named special case, not a restriction of the general algorithm request. The research summary reports algorithms in some cases but supplies no precise references or scope for them.
\subsection{Short English statement}
Find an exact algorithm for the Veech group of an arbitrary translation surface or quadratic differential, beyond special cases such as square-tiled surfaces. Specify the input encoding and what description of the possibly infinite group must be returned.
\subsection{Sources}
\begin{itemize}
\item Local UnsolvedMath registry, permanent ID 11000214, source code catalogue entry; no proposer or proposal year supplied.
\item Farb, edited volume \emph{Problems on Mapping Class Groups and Related Topics}, Problem 7, PDF page 259; \url{https://www.math.uchicago.edu/~farb/papers/mcgbook.pdf}. The registry records PDF-font-structure extraction and access on 2026-07-29; the inventory distinguishes 329 labeled headings by source page.
\item Source labels [Zo3] and [HuLe1] for square-tiled surfaces; full bibliography entries are not reproduced. The research summary supplies no individual citations for later algorithms.
\end{itemize}
\end{document}
